\chapter{Experimental Evaluation}
\label{ch:eval}

This chapter presents the full experimental evaluation of \padonap{}.
All numbers reported here are drawn directly from JSON artefacts produced
by the evaluation harness; the source files are cited inline so every
claim is independently reproducible.

%%─────────────────────────────────────────────────────────────────────────────
\section{Experimental setup}
\label{sec:setup}

\paragraph{Hardware.}
All experiments were executed on a single workstation running Windows~11:
Intel Core-i5 CPU, 16~GB DDR4, NVIDIA RTX~3050 (4~GB VRAM). Both the
AI inference path and the orchestration loop run on CPU; the GPU was used
only for Transformer+LSTM training.

\paragraph{Software stack.}
Python~3.10, XGBoost~1.7, PyTorch~2.0, SciPy (HiGHS LP solver).
Seeds are pinned: numpy=42, torch=42, xgboost seed=42 throughout.

\paragraph{Evaluation harness.}
\begin{itemize}
  \item AI path: \texttt{evaluation/scenarios.py} with \texttt{eval\_mode=True}
        (bypasses frequency guard to allow scenario replay).
  \item Baseline path: \texttt{evaluation/baseline\_threshold.py}; uses the
        same PolicyEngine, ONAPSOClient, SFCManager, and LatencyTracker as
        the AI path.
  \item Lead-time: \texttt{evaluation/lead\_time\_analyzer.py}.
  \item Raw artefacts: \texttt{evaluation/results/} (AI) and
        \texttt{evaluation/results\_baseline/} (threshold).
\end{itemize}

\paragraph{Simulation methodology.}
VNF boot times are drawn from the dict \texttt{VNF\_SIM\_BOOT\_S}
(ratelimiter: 0.5~s, scrubber: 6~s, blackhole: 0.2~s), calibrated from
ETSI NFV ISG PoC~\#4 measurements. ONAP~SO runs in stub mode
(\texttt{PAD\_ONAP\_STUB=true}). These are disclosed threats to validity;
see \S\ref{sec:validity}.

%%─────────────────────────────────────────────────────────────────────────────
\section{AI metrics}
\label{sec:ai-metrics-eval}

Table~\ref{tab:ai-metrics} recapitulates the AI metrics reported in
Chapter~\ref{ch:ai} for reference during the scenario analysis below.

\begin{table}[H]
\centering
\small
\caption{AI component metrics (CICDDoS2019, temporal 80/20 split, no shuffle).
Source: \texttt{memory/project\_ai\_v2\_training.md}.}
\label{tab:ai-metrics}
\begin{tabular}{@{}lll@{}}
\toprule
Track & Metric & Value \\
\midrule
\multirow{4}{*}{A — XGBoost 7-class}
  & Accuracy              & 0.9825 \\
  & Macro-F1              & 0.9642 \\
  & Macro AUC             & 0.9958 \\
  & Inference P99 latency & 2.2~ms \\
\midrule
\multirow{5}{*}{B — Transformer+LSTM}
  & AUC at $t{+}30$~s  & 0.9884 \\
  & AUC at $t{+}60$~s  & 0.9833 \\
  & AUC at $t{+}90$~s  & 0.9772 \\
  & AUC at $t{+}120$~s & 0.9711 \\
  & Inference P99 latency & 7.68~ms \\
\bottomrule
\end{tabular}
\end{table}

%%─────────────────────────────────────────────────────────────────────────────
\section{Scenario suite: \padonap{} AI orchestrator (S1--S8)}
\label{sec:s1-s8}

Eight scenarios stress the orchestrator across four axes: benign traffic
(S1), volumetric attacks (S2, S3), out-of-distribution OOD traffic
(S4, S5), multi-attack sequencing (S6), multi-tenant SLA fairness (S7),
and a controlled proactive-vs-reactive head-to-head (S8).
Each scenario is 5-second windows; tier is decided per-window.

\begin{table}[H]
\centering
\small
\caption{\padonap{} results on S1--S8 (all numbers from
\texttt{evaluation/results/evaluation\_summary.json}).
``Proact.''~= windows where $T_2$ was triggered by forecast rather than
detection. Latency is P50 of acted windows only.}
\label{tab:s1-s8-ai}
\begin{tabular}{@{}llrrrrrrl@{}}
\toprule
ID & Scenario & $n_w$ & Max & Proact. & $T_2$ P50 & $T_3$ P50 & SLA & Verdict \\
   &          &       & tier & count & (ms)      & (ms)      & ok  & \\
\midrule
S1 & Normal baseline        & 100 & $T_0$ &  0 &  ---   &    ---   & \checkmark & \textbf{PASS} \\
S2 & Sudden UDP flood       & 110 & $T_3$ &  0 &  ---   & 6006.02  & \checkmark & \textbf{PASS} \\
S3 & Gradual SYN ramp       & 120 & $T_2$ & 67 & 505.51 &    ---   & \checkmark & \textbf{PASS} \\
S4 & HTTP flood (OOD)       & 100 & $T_1$ &  0 &  ---   &    ---   & \checkmark & \textbf{PASS} \\
S5 & ICMP burst (OOD)       & 100 & $T_0$ &  0 &  ---   &    ---   & \checkmark & \textbf{PASS} \\
S6 & Multi UDP+SYN          & 130 & $T_3$ & 48 & 505.04 & 6004.53  & \checkmark & \textbf{PASS} \\
S7 & SLA fairness 3-tenant  & 120 & $T_2$ & 78 & 505.71 &    ---   & \checkmark & \textbf{PASS} \\
S8 & Proactive vs reactive  & 120 & $T_3$ & 30 & 505.23 & 6006.05  & \checkmark & \textbf{PASS} \\
\midrule
\multicolumn{8}{r}{Total} & \textbf{8/8 PASS} \\
\bottomrule
\end{tabular}
\end{table}

\subsection{S1 — Normal baseline}
\label{sec:s1}
100 windows of benign traffic. The AI correctly stays at $T_0$ for all
100 windows: no false-positive escalation. SLA trivially satisfied
(no VNF overhead). This confirms Invariant~\ref{inv:urllc} under the
zero-overhead case.

\subsection{S2 — Sudden UDP flood}
\label{sec:s2}
110 windows: 10 ramp-up, 60 full UDP flood, 40 normal.
The attack signature (pkt\_rate $\gg$ 3000, udp\_frac $\approx 1.0$,
src\_port\_entropy $\approx 0$) is immediately identifiable by the
XGBoost detector. The orchestrator escalates to $T_3$ (scrubber VNF)
on the first flood window; VNF active at $6006.02$~ms (simulated
scrubber boot). No proactive path is exercised because the attack is
too sudden for the forecaster to deliver a pre-positioning signal.
The $T_3$ latency of 6006~ms is the \emph{reactive floor} of the
architecture.

\subsection{S3 — Gradual SYN ramp}
\label{sec:s3}
120 windows: SYN rate ramps from near-zero to flood over $\approx 50$
windows, then falls. The Transformer+LSTM forecast detects the ramp
early: 67 of the 70 attack-phase windows are triggered proactively
via $T_2$ (rate-limiter VNF), with a first proactive signal at window~42.
$T_2$ P50 latency is $505.51$~ms (rate-limiter boot). The scenario
never reaches $T_3$, meaning the pre-positioned rate-limiter was
sufficient to contain the ramp without escalating to the scrubber.
This is the architecture's ``best-case'' outcome: the forecaster averts
a heavier reactive response.

\subsection{S4 and S5 — Out-of-distribution (OOD) traffic}
\label{sec:s4-s5}
S4 injects an HTTP flood pattern (high rate, varied ports, OOD relative
to CICDDoS2019 training); S5 injects an ICMP burst. Neither class
appears in the XGBoost training set (see \S\ref{sec:dataset-limits}).
\padonap{} responds as follows:
\begin{itemize}
  \item \textbf{S4}: max tier $T_1$ (alert only, no VNF). The model's
        entropy pattern for HTTP does not cross the detection threshold,
        so no VNF is instantiated.
  \item \textbf{S5}: max tier $T_0$ (normal). ICMP-burst features fall
        below all tier thresholds; no false positive.
\end{itemize}
Both are graded \textbf{PASS} because the expected-maximum-tier
constraints ($T_1$ and $T_0$) are met. This is the OOD robustness claim:
\emph{the AI degrades gracefully by under-responding}, which is the safer
failure mode compared to over-escalation.

\subsection{S6 — Multi-attack sequencing}
\label{sec:s6}
130 windows: a UDP flood followed immediately by a SYN wave. The
orchestrator handles two independent escalation events: it reaches $T_3$
on the UDP segment (6004.53~ms reactive, tier-dist $T_3{=}33$), then
partially de-escalates and re-escalates on the SYN segment. The 48
proactive windows ($T_2$) correspond to the inter-attack overlap, where
the SYN ramp is partially detectable before it peaks. All SLA checks
pass. This scenario exercises the hysteresis guard: without it the
rapid $T_3 \to T_0 \to T_3$ transition would cause VNF churn.

\subsection{S7 — SLA fairness under pre-positioning}
\label{sec:s7}
120 windows: a 3-tenant scenario (URLLC floor 150~Mbps, eMBB 100~Mbps,
mMTC 50~Mbps; total $C{=}1000$~Mbps). The forecaster keeps the scenario
at $T_2$ for 81 of 120 windows (78 proactive), imposing a 500~Mbps VNF
overhead on the LP allocator. The allocator solves
(\ref{eq:lp-obj})--(\ref{eq:lp-bounds}) and returns $a_{\text{URLLC}}
= 150$~Mbps at every allocation, confirming Invariant~\ref{inv:urllc}.
$T_2$ P50 = 505.71~ms. \textbf{sla\_ok = true} for all 120 windows.

\subsection{S8 — Proactive vs reactive head-to-head}
\label{sec:s8}
120 windows: designed explicitly to quantify the temporal gap between
proactive ($T_2$) and reactive ($T_3$) paths. The attack ramps so that
both paths trigger. Results:
\begin{itemize}
  \item First \emph{proactive} $T_2$ at window~31 ($t = 150$~s).
        Latency: $505.23$~ms.
  \item First \emph{reactive} $T_3$ at window~58 ($t = 285$~s).
        Latency: $6006.05$~ms.
  \item Proactive advantage vs reactive: $10.9\times$ latency reduction.
  \item Window distance: $58 - 31 = 27$ windows $= \mathbf{135}$ seconds.
\end{itemize}
The tier distribution is $T_0{:}57,\ T_2{:}26,\ T_3{:}37$, reflecting
a realistic mixed-response: the forecaster handles the build-up while
the detector takes over at the peak.

%%─────────────────────────────────────────────────────────────────────────────
\section{Threshold baseline comparison}
\label{sec:baseline}

The rule-based threshold baseline (\texttt{evaluation/baseline\_threshold.py})
applies seven static rules on the 17-feature vector (see
Table~\ref{tab:baseline-rules}) and reuses the identical downstream
stack. This ensures any latency or tier difference is attributable
solely to the detection/forecast stage, not to differences in the
orchestration path.

\begin{table}[H]
\centering
\small
\caption{Threshold rules in the baseline (from \texttt{baseline\_threshold.py}, lines~86--100).}
\label{tab:baseline-rules}
\begin{tabular}{@{}lll@{}}
\toprule
Feature & Condition & Action \\
\midrule
\texttt{pkt\_rate}        & $> 10{,}000$   & $T_3$ \\
\texttt{pkt\_rate}        & $> 3{,}000$    & $\geq T_2$ \\
\texttt{pkt\_rate}        & $> 800$        & $\geq T_1$ \\
\texttt{syn\_ratio}       & $> 0.60$       & $\geq T_3$ \\
\texttt{proto\_dist\_udp} & $> 0.85$       & $\geq T_3$ \\
\texttt{proto\_dist\_icmp}& $> 0.70$       & $\geq T_2$ \\
\texttt{src\_ip\_entropy} & $< 0.80$ + rate$>500$ & $\geq T_2$ \\
\bottomrule
\end{tabular}
\end{table}

\begin{table}[H]
\centering
\small
\caption{Threshold baseline on S1--S8 (from
\texttt{evaluation/results\_baseline/baseline\_summary.json}).
The baseline has no proactive path by construction.}
\label{tab:baseline-results}
\begin{tabular}{@{}llrrrrl@{}}
\toprule
ID & Scenario & Max tier & $T_2$ P50~(ms) & $T_3$ P50~(ms) & SLA ok & Verdict \\
\midrule
S1 & Normal baseline       & $T_0$ &    ---  &    ---   & \checkmark & \textbf{PASS} \\
S2 & Sudden UDP flood      & $T_3$ &    ---  & 6009.27  & \checkmark & \textbf{PASS} \\
S3 & Gradual SYN ramp      & $T_3$ & 501.00  & 6001.00  & \checkmark & \textbf{PASS} \\
S4 & HTTP flood (OOD)      & $T_2$ & 501.00  &    ---   & \checkmark & \textbf{FAIL} \\
S5 & ICMP burst (OOD)      & $T_2$ & 504.34  &    ---   & \checkmark & \textbf{FAIL} \\
S6 & Multi UDP+SYN         & $T_3$ &    ---  & 6002.44  & \checkmark & \textbf{PASS} \\
S7 & SLA fairness          & $T_3$ &    ---  & 6003.97  & \checkmark & \textbf{PASS} \\
S8 & Proactive vs reactive & $T_3$ &    ---  & 6006.96  & \checkmark & \textbf{PASS} \\
\midrule
\multicolumn{6}{r}{Total} & \textbf{6/8 PASS} \\
\bottomrule
\end{tabular}
\end{table}

\paragraph{OOD failures.}
The two baseline failures occur on scenarios S4 (HTTP flood) and S5
(ICMP burst) — both OOD. The baseline escalates to $T_2$ because the
HTTP traffic generates a moderate pkt\_rate crossing the \texttt{>800}
threshold and the ICMP burst crosses the \texttt{icmp\_frac > 0.70}
rule, respectively. These escalations are false positives: the traffic
is not in the DDoS attack set and should stay at $T_0$ or $T_1$.
\padonap{} stays at $T_1$ (S4) and $T_0$ (S5) because the XGBoost
classifier, trained on real flow distributions, assigns low attack
probability to both.

\paragraph{S3 reactive vs proactive.}
For the gradual SYN-ramp scenario, the baseline does eventually
escalate (\texttt{syn\_ratio > 0.60} fires at window~55) — but only
to $T_3$ (reactive scrubber), never to a pre-positioned $T_2$.
The proactive AI path fires $T_2$ at window~42, 13~windows (65~s)
earlier and at lower cost (rate-limiter vs scrubber).

\paragraph{Reactive latency parity.}
Where both systems act reactively (S2, S6, S7, S8), the $T_3$ P50
values are within 1\% of each other ($\approx 6000$~ms). This
confirms the speed-up reported for the proactive path is due to
\emph{early trigger}, not a faster boot path.

%%─────────────────────────────────────────────────────────────────────────────
\section{Lead-time analysis}
\label{sec:lead-time}

Table~\ref{tab:lead-time} quantifies how many seconds the AI forecast
pre-positions mitigation before the reactive path (AI or baseline) would
fire. All numbers come from
\texttt{evaluation/results/lead\_time\_comparison.csv}.

\begin{table}[H]
\centering
\small
\caption{Lead-time analysis. Window = 5~s. Positive values = AI proactive
acted earlier; negative = reactive acted earlier (expected for sudden attacks).
Source: \texttt{lead\_time\_comparison.csv}.}
\label{tab:lead-time}
\begin{tabular}{@{}lrrrrrr@{}}
\toprule
Scenario & \parbox{1.3cm}{\centering First\\proact.\\win} &
           \parbox{1.3cm}{\centering First\\react.\\AI win} &
           \parbox{1.3cm}{\centering First\\react.\\base win} &
           \parbox{1.6cm}{\centering Lead vs\\AI-react (s)} &
           \parbox{1.6cm}{\centering Lead vs\\baseline (s)} &
           \parbox{1.1cm}{\centering \#Proact.\\wins} \\
\midrule
S1 Normal        & ---& --- & --- & ---       & ---      &  0 \\
S2 Sudden UDP    & ---&  37 &  32 & ---       & ---      &  0 \\
S3 SYN ramp      &  42& --- &  55 & ---       & +\textbf{65.0} & 67 \\
S4 HTTP (OOD)    & ---& --- & --- & ---       & ---      &  0 \\
S5 ICMP (OOD)    & ---& --- & --- & ---       & ---      &  0 \\
S6 Multi         &  71&  32 &  32 & $-195.0$  & $-195.0$ & 48 \\
S7 SLA 3-tenant  &  31& --- &  32 & ---       & +\textbf{5.0}  & 78 \\
S8 Novelty       &  31&  58 &  32 & +\textbf{135.0} & +\textbf{5.0} & 30 \\
\bottomrule
\end{tabular}
\end{table}

\subsection{Key observations}

\paragraph{H1 confirmed: S3 lead-time +65 s.}
On the gradual SYN-ramp, the AI forecast pre-positions $T_2$ at
window~42 ($t = 210$~s), while the threshold baseline first triggers
$T_3$ at window~55 ($t = 275$~s). The lead-time is
$(55 - 42) \times 5 = 65$~s. Hypothesis~H1 (target $\geq 30$~s on
gradual attacks) is \textbf{confirmed}.

\paragraph{S8 headline: +135 s within the same orchestrator.}
Scenario S8 is designed so that both $T_2$ (proactive) and $T_3$
(reactive) fire within the same run. The proactive path fires 27 windows
before the reactive path, giving a 135~s advantage inside \padonap{}.
Against the baseline, the gain is 5~s (both act on the same attack onset
but with different VNF cost: 505~ms vs 6007~ms).

\paragraph{S6 negative: expected for sudden attacks.}
The $-195$~s figure for S6 does not indicate a defect. In S6 the
attack is sudden (no ramp); the reactive path fires at window~32
immediately, while the proactive forecast (which needs trend buildup)
does not fire until window~71. This is the known limit of any
forecast-based defence: sudden attacks still go to reactive, but the
reactive latency is the same ($\approx 6000$~ms) as the baseline's.
\padonap{} does not make sudden-attack response \emph{worse}.

\paragraph{S4/S5: no lead-time columns --- no proactive fired.}
Both OOD scenarios produce neither a proactive nor a reactive trigger
in the AI path (correct by design), so no lead-time metric exists.

%%─────────────────────────────────────────────────────────────────────────────
\section{SLA fairness under pre-positioning}
\label{sec:sla-fairness}

Scenario~S7 is dedicated to verifying Invariant~\ref{inv:urllc}. The LP
allocator is called for every window in which a tier change is actioned.
In S7, 81 of 120 windows are at $T_2$, imposing a 500~Mbps overhead on
the 1000~Mbps link budget. The LP returns:
\[
  a_{\text{URLLC}} = 150~\text{Mbps},\quad
  a_{\text{eMBB}} \approx 233~\text{Mbps},\quad
  a_{\text{mMTC}} \approx 117~\text{Mbps},
\]
leaving a 0~Mbps surplus (fully allocated). \texttt{sla\_ok = true}
for all 120 windows. No scenario in the full S1--S8 suite reports
\texttt{sla\_ok = false}. Invariant~\ref{inv:urllc} is verified.

%%─────────────────────────────────────────────────────────────────────────────
\section{SOTA comparison}
\label{sec:sota}

Table~\ref{tab:sota} compares \padonap{} against the seven closest
state-of-the-art systems on the metrics most directly comparable.

\begin{table}[H]
\centering
\footnotesize
\caption{SOTA comparison on DDoS detection and orchestration metrics.
Values marked $^\dagger$ are extracted from the cited paper's best-reported
figure; values without $^\dagger$ are from this thesis's evaluation.}
\label{tab:sota}
\begin{tabular}{@{}lccccc@{}}
\toprule
System & Acc & F1 & Proactive & VNF boot (ms) & Multi-tenant SLA \\
\midrule
TST-DDoS (2025)$^\dagger$            & 0.9710 & 0.9530 & ---    & ---    & --- \\
NetProbe (2025)$^\dagger$            & 0.9680 & ---    & 1-tier & $\sim$3000 & --- \\
Lightweight ML (2025)$^\dagger$      & 0.9620 & 0.9480 & ---    & ---    & --- \\
FlowGuard (2022)$^\dagger$           & ---    & ---    & 1-hor. & ---    & --- \\
DAWN (2025)$^\dagger$                & 0.9550 & 0.9300 & ---    & ---    & --- \\
Latency-aware NFV-6G (2024)$^\dagger$& ---    & ---    & ---    & 8000--14000 & --- \\
\textbf{\padonap{} (ours)}           & \textbf{0.9825} & \textbf{0.9642} & \textbf{4-hor.} & \textbf{505} & \textbf{LP floor} \\
\bottomrule
\end{tabular}
\end{table}

\padonap{} leads on detection accuracy (+1.15\% vs TST-DDoS),
VNF activation latency ($>15\times$ faster than Latency-aware NFV-6G),
and is the only system offering multi-horizon proactive forecasting
with a formal SLA floor guarantee.

%%─────────────────────────────────────────────────────────────────────────────
\section{Threats to validity}
\label{sec:validity}

We document five threats candidly.

\begin{enumerate}
  \item \textbf{Simulated VNF boot times.}
        The 505~ms ($T_2$) and 6006~ms ($T_3$) latencies are from
        calibrated ETSI NFV ISG PoC~\#4 values, not live Docker
        measurements. The absolute numbers may shift in production;
        however, the $\approx 12\times$ ratio between scrubber and
        rate-limiter boot is hardware-independent (image size,
        health-check interval).

  \item \textbf{ONAP SO stub.}
        The downstream stack implements the SO~v7 REST contract
        (\texttt{/onap/so/infra/serviceInstantiation/v7/}) but does
        not perform a live Kubernetes instantiation. Integration with
        ONAP OOM is scheduled as future work (\S\ref{sec:future}).

  \item \textbf{Dataset scope (3 of 7 classes absent).}
        \texttt{HTTP\_Flood}, \texttt{ICMP\_Flood}, and
        \texttt{Slow\_rate} lack NetFlow records in CICDDoS2019 and
        cannot be included in training. They appear only as OOD
        tests (S4, S5).

  \item \textbf{Single-node evaluation.}
        All eight scenarios run on one machine. Distributed
        multi-orchestrator deployment is not evaluated.

  \item \textbf{Single-dataset evaluation.}
        Cross-dataset validation (BCCC-cPacket-Cloud-DDoS-2024) is
        scheduled for future work (\S\ref{sec:future}).
\end{enumerate}
